\section{Bargaining about another player's opportunities}
\label{sec:bargaining}

\subsection{Responses and the composition of offers}
Across 526 negotiation threads, players accept 113 offers, counter 394 and reject 413. Acceptance is legal in 918 of the 920 response opportunities. Sol and Opus mostly counter, in 140/235 and 141/189 responses. Gemini and Grok mostly reject, in 180/242 and 149/254. Each nevertheless accepts between 26 and 30 offers. Those similar settlement totals emerge from different ways of responding to disagreement.

These are applied responses. Seven of Grok's 28 acceptances conflict with the action named in its assistant content. The distinction between a submitted call and accompanying text therefore applies on both sides of a transaction, not only to the offers initiating it.

Table~\ref{tab:opening-all} reports every opening-offer form and completion category. Among swaps, offers completing only the proposer receive 95/116 rejections, compared with 37/72 completing neither participant. Both-completing swaps receive 34/48 rejections. Pure cash bids give a different estimate: 42/70 proposer-only bids and 67/103 neither-completing bids are rejected. The latter difference is negative in this cohort, but its uncertainty does not establish a zero or negative population effect. Later swaps show much smaller differences, 8/39 versus 7/34. They are selected by the preceding negotiation.

\begin{table*}[t]\centering\small
\begin{tabular}{@{}lrrrr@{}}\toprule
Newly completes & Swap & Cash bid & Cash ask & Other\\\midrule
Proposer only & 116 [1/20/95] & 70 [2/26/42] & 0 [0/0/0] & 1 [0/0/1] \\
Recipient only & 6 [1/3/2] & 0 [0/0/0] & 41 [1/14/26] & 4 [0/1/3] \\
Both & 48 [3/11/34] & 0 [0/0/0] & 0 [0/0/0] & 0 [0/0/0] \\
Neither & 72 [17/18/37] & 103 [4/32/67] & 55 [3/12/40] & 8 [0/0/8] \\
\bottomrule\end{tabular}
\caption{Every opening offer where acceptance is legal: N [accept/counter/reject]. Swaps transfer deeds in both directions; cash bids request deeds; cash asks offer deeds. The other category retains bundles outside these definitions. Completion concerns proposed ownership, not immediate ability to build. The appendix gives all later offers too.}\label{tab:opening-all}
\end{table*}


Proposer composition matters to the opening-swap difference. Grok writes 36 of the 72 neither-completing swaps and fifteen of their seventeen acceptances. Thirteen of those accepted offers originate in calls whose assistant content names a different action. Table~\ref{tab:proposer} shows that each observed proposer stratum has a positive rejection difference, but the comparison cells vary greatly in size. Removing Grok-authored offers changes the pooled difference from 30.5 to 16.7 percentage points. Removing only threads opened by conflicting calls instead gives 93/112 versus 26/43, a 22.6-point difference. These exclusions describe different populations; neither recreates play under a corrected interface.

\begin{table}[t]\centering\small
\begin{tabular}{@{}lrrr@{}}\toprule
Proposer & Own only & Neither & Gap, pp\\\midrule
Sol & 44/50 & 9/11 & 6.2 \\
Gemini & 4/8 & 0/2 & --- \\
Opus & 42/50 & 15/23 & 18.8 \\
Grok & 5/8 & 13/36 & 26.4 \\
\bottomrule

\end{tabular}
\caption{Opening swaps, rejected/received. ``Own only'' refers to the proposer. A dash omits a gap when either cell has fewer than five offers. Compensation, deeds and positions differ within rows.}\label{tab:proposer}
\end{table}

Across the five seed blocks, the unfiltered opening-swap difference stays positive. Without Grok authors, one block reverses to $-23.5$ points; the five comparison cells contain only 10, 11, 8, 4 and 3 offers. Leaving one block out gives 28.7--33.3 points overall and 9.4--22.1 without Grok authors. Figure~\ref{fig:stability} shows the block values and omission ranges; Appendix~\ref{app:comparisons} prints every block's counts. Resampling schemes that also vary the representation of rotations give wider ranges than intact-block resampling. The no-Grok range includes zero under each of the three reported schemes.

\begin{figure*}[t]\centering
\includegraphics[width=.87\textwidth]{bargaining_stability.pdf}
\caption{Opening rejection differences between proposer-only and neither completion. Blue circles show pooled differences, orange ticks the five seed-block differences, and gray segments the range after leaving out each block. Rows use different populations. The no-Grok comparison cells contain only 10, 11, 8, 4 and 3 offers; Table~\ref{tab:block-counts} gives all block counts.}\label{fig:stability}
\end{figure*}

Recipient composition matters too. The pooled cash-bid difference is $-5.0$ points, while the four within-recipient differences range from $-0.2$ to $+8.5$. Appendix~\ref{app:comparisons} gives their counts. The forms also have different proposers: Opus writes 110/173 focal opening cash bids, including 40/70 proposer-only and 70/103 neither-completing bids; 105 use action-only recovery. Cash and property compensation and recipients' portfolios differ across offers. Board list prices omit complementarity and blocking value. The census establishes the pattern of responses; the statements below supply direct evidence of the considerations agents express.

\subsection{Explicit consideration of rival development}
The separate semantic sample consists of seventy-four rejections of offers completing the recipient's group for which acceptance was legal. The two Astra readings agree that 61 name a trade-linked consequence for the rival: 27/38 recipient-only and 34/36 both-completing offers. These populations are disjoint from the proposer-only/neither comparison above. They document expressed consideration of the rival, not an explanation for the swap/cash difference. Gemini and Grok contribute 56 of the agreed cases; G1 and G15 contribute 31.

\TranscriptTitle{g15}{Rejecting one's own completion while citing the rival's}
At G15 turn 63, Sol offers two greens for Ventnor and \$100, completing both groups. Grok has \$333 and Sol \$730. Grok's note at turn 33 already rules out selling Ventnor back to Sol, and its notes at turns 61 and 62 repeat the blocking policy. Immediately before the offers, Opus publicly urges the others not to give Sol a final completing deed; Grok receives that appeal. Its refusal follows the earlier private policy as well as this public warning:
\TranscriptQuote{fx-g15-refusal-note}
Sol subsequently offers the greens and Kentucky, then adds \$250 and explicitly compares Grok's resulting \$583 with its own \$480. Grok still refuses:
\TranscriptQuote{fx-g15-final-refusal}
The first offer leaves the stated \$830 and \$233. A full first layer on green costs \$600, although one legal house costs \$200. The objection concerns the rival's development, even when Grok would receive its own monopoly. All three offers fail. Eighteen turns later Grok's blockers pass to Gemini through bankruptcy; Sol never completes either group. This continuation reveals the stakes of withholding without identifying whether refusal improved Grok's chances.
\TranscriptSource{G15 d168--175, standard responses. The quoted public message occurs in the next counterparty request.}

\subsection{Prices and counteroffers express different choices}
Countering often retains the disputed exchange. In 111 counters to proposer-only completion offers, 79 preserve the entire noncash exchange while improving the counterer's cash. The corresponding figure for non-completing deed requests is 81/110. Repricing is common in both populations. Among unchanged-bundle concessions improving the recipient's cash, 9/58 proposer-only completion offers and 17/56 neither-completing deed requests are accepted. More money often leaves disagreement unresolved; these selected sequences do not identify a price response.

Median single-deed asks relative to list price are higher for completing buyers in all seven seller/mortgage-status cells containing both buyer types; the completing-buyer cells contain 1--17 asks. More closely overlapping histories are mixed. Six same-game, same-seller, same-deed comparisons retain mortgage status: completing-buyer asks are higher in four, lower in one and equal in one. The seventh changes mortgage status and its ordering depends on the summary used. Time, buyer and negotiation history still differ. Appendix~\ref{app:comparisons} gives every pair.

\TranscriptTitle{g2}{Pricing the buyer's development}
Opus holds two yellows and bids \$325 for Sol's Marvin Gardens. Sol responds:
\TranscriptQuote{g2-d124-price}
One turn earlier Grok, which owned no yellow, bid \$200 for the deed as a blocker. Sol asked \$450, and Grok replied that this was a price for a completing buyer. Grok then publicly warned Sol against completing Opus's yellows; the warning appears in Sol's next prompt. Sol had already recorded the blocking policy four turns earlier and publicly priced its control of yellow before the warning. Sol's note discusses Opus's remaining \$520 after the lower bid, but incorrectly calls it enough for nine yellow houses; it buys three. Opus calls \$450 final while privately retaining a \$475 ceiling. Sol refuses:
\TranscriptQuote{g2-d128-keep}
At turn 69 the exchange reverses direction. Opus has \$17 and sells its two yellows to Sol for two railroads and \$100. Sol builds, later liquidates those buildings in the G2 d216 episode of Table~\ref{tab:necessary}, collects \$44 undeveloped rent and no developed yellow rent, then loses the group to Gemini through bankruptcy. The sequence connects an attempt to buy a blocker, buyer-specific asking prices and a later change in who can finance the group.
\TranscriptSource{G2 d119--128 and d197--198. Opus's first exchange uses standard responses; its later sale uses recovery.}

\TranscriptTitle{g19}{Accepting terms while correcting their justification}
Grok proposes Virginia and \$150 for States:
\TranscriptQuote{g19_d135_false_public}
Both deeds are pink. The same response's note says to abort and its assistant content names an end action, but the received proposal reaches Sol:
\TranscriptQuote{g19_d136_correction_accept}
The exchange executes and completes no group. Sol explicitly separates the false premise from acceptable terms. Grok's conflicting channels make the case unsuitable as evidence of one coherent deceptive plan, but the recipient's correction and acceptance are directly observed.
\TranscriptSource{G19 d135--136, standard responses. Sol's stored prompt contains the quoted false message.}
